\subsection{Analytic Semigroups}\label{sec:analytic-semigroups}

\subsubsection{Analytic Semigroups and Their Generators}\label{subsec:analytic-semigroups-generators}

We shall now discuss a class of semigroups which in many respects allows for much stronger results than strongly continuous semigroups. In particular we shall obtain:\begin{enumerate}\item Better regularity of solutions to initial-value problems.\item Better results concerning perturbations of the infinitesimal generator.\item A relationship between the type of the semigroup and the spectrum of the infinitesimal generator.\end{enumerate}We begin with the definition of an analytic semigroup.

\begin{Definition}
\label{def:analytic-semigroup}
A strongly continuous semigroup $\exp(At)$ is called an \textbf{analytic semigroup} if the following conditions hold:\begin{enumerate}\item For some $\theta\in(0,\pi/2)$, $\exp(At)\in\mathcal{L}(X)$ can be extended to $t\in\Delta_\theta=\{0\}\cup\{t\in\mathbb{C}\mid|\arg t|<\theta\}$ and the conditions of Definition~\ref{def:strongly-continuous-semigroup} hold for $t\in\Delta_\theta$.\item For $t\in\Delta_\theta\setminus\{0\}$, $\exp(At)$ is analytic in $t$ (in the sense of the uniform operator topology).\end{enumerate}
\end{Definition}

The most important question now is how to characterize the infinitesimal generators of analytic semigroups. This is addressed by the following theorem.

\begin{Theorem}
\label{thm:analytic-generator-sectorial-resolvent}
A closed, densely defined operator $A$ in $X$ is the generator of an analytic semigroup if and only if there exists $\omega\in\mathbb{R}$ such that the half-plane $\operatorname{Re}\lambda>\omega$ is contained in the resolvent set of $A$ and, moreover, there is a constant $C$ such that
\begin{equation}\label{eq:analytic-semigroup-half-plane-resolvent-estimate}
\|R_\lambda(A)\|\leq C/|\lambda-\omega|
\end{equation}
for $\operatorname{Re}\lambda>\omega$. If this is the case, then actually the resolvent set contains a sector $|\arg(\lambda-\omega)|<\frac{\pi}{2}+\delta$ for some $\delta>0$, and an analogous resolvent estimate holds in this sector. Moreover, the semigroup is represented by
\begin{equation}\label{eq:analytic-semigroup-contour-integral-representation}
e^{At}=\dfrac{1}{2\pi i}\int_\Gamma e^{\lambda t}(\lambda I-A)^{-1}\,\mathrm{d}\lambda,
\end{equation}
where $\Gamma$ is any curve from $e^{-i\phi}\infty$ to $e^{i\phi}\infty$ such that $\Gamma$ lies entirely in the set $\{|\arg(\lambda-\omega)|\leq\phi\}$. Here $\phi$ is any angle such that $\frac{\pi}{2}<\phi<\frac{\pi}{2}+\delta$.
\end{Theorem}

\begingroup
\renewcommand{\endProof}{\par}
\begin{Proof}
The necessity of the condition is not hard to see. Assume $\exp(At)$ is an analytic semigroup defined in $\Delta_\theta$. It is easy to show that for any $\psi<\theta$, there are constants $M$ and $\omega$ such that $\|\exp(At)\|\leq M\exp(\omega|t|)$ for $t\in\Delta_\psi$; the proof proceeds along the lines of Theorem~\ref{thm:semigroup-logarithmic-growth-type}. From the proof of the Hille-Yosida theorem, we know that
\begin{equation}\label{eq:analytic-generator-necessity-laplace-resolvent}
-R_\lambda(A)=\int_0^\infty e^{-\lambda t}T(t)\,\mathrm{d}t
\end{equation}
for $\operatorname{Re}\lambda>\omega$. Without loss of generality, let $\operatorname{Im}\lambda>0$. Then we shift the contour to the ray $\arg t=-\delta$, $\delta=\max(\psi,\arg\lambda)$ and obtain
\begin{equation}\label{eq:analytic-generator-rotated-laplace-integral}
-R_\lambda(A)=\int_0^\infty\exp(-\lambda e^{-i\delta}t)T(e^{-i\delta}t)e^{-i\delta}\,\mathrm{d}t,
\end{equation}
and hence
\begin{equation}\label{eq:analytic-generator-rotated-resolvent-bound}
\|R_\lambda(A)\|\leq\dfrac{M}{\operatorname{Re}(\lambda\exp(-i\delta))-\omega}\leq\dfrac{C}{|\lambda-\omega|}.
\end{equation}

\end{Proof}
\endgroup

Assume now that \eqref{eq:analytic-semigroup-half-plane-resolvent-estimate} holds for $\operatorname{Re}\lambda>\omega$. If $\lambda-\omega=Re^{i\theta}$, $\theta>\pi/2$, and $\phi$ is any angle less than $\pi/2$, we can write
\begin{equation}\label{eq:analytic-generator-sector-resolvent-extension-splitting}
A-\lambda I=A-(\omega+Re^{i\phi})I-R(e^{i\theta}-e^{i\phi})I,
\end{equation}
and the inverse can be represented as a geometric series as long as $|e^{i\theta}-e^{i\phi}|\leq C$. This shows that the resolvent set extends to a sector beyond the half-plane $\operatorname{Re}\lambda>\omega$, and a resolvent estimate is obtained by term-by-term estimation of the geometric series. In particular, this makes \eqref{eq:analytic-semigroup-contour-integral-representation} well defined.

We next show that \eqref{eq:analytic-semigroup-contour-integral-representation} does indeed represent an analytic semigroup generated by $A$. Clearly, \eqref{eq:analytic-semigroup-contour-integral-representation} is defined for $|\arg t|<\phi-\frac{\pi}{2}$ and it is analytic in $t$. We next check the semigroup property. We denote by $T(t)$ the right-hand side of \eqref{eq:analytic-semigroup-contour-integral-representation} and we take $\Gamma^{\prime}$ to be a path with the same properties as $\Gamma$, but lying entirely to the right of $\Gamma$. Then we have\[\begin{aligned}
T(s)T(t)&=\left(\dfrac{1}{2\pi i}\right)^2\int_{\Gamma^{\prime}}\int_\Gamma e^{\lambda^{\prime}s+\lambda t}(\lambda^{\prime}I-A)^{-1}(\lambda I-A)^{-1}\,\mathrm{d}\lambda\,\mathrm{d}\lambda^{\prime}\\
&=\left(\dfrac{1}{2\pi i}\right)^2\Bigg[\int_{\Gamma^{\prime}}e^{\lambda^{\prime}s}(\lambda^{\prime}I-A)^{-1}\int_\Gamma e^{\lambda t}(\lambda-\lambda^{\prime})^{-1}\,\mathrm{d}\lambda\,\mathrm{d}\lambda^{\prime}\\
&\qquad-\int_\Gamma e^{\lambda t}(\lambda I-A)^{-1}\int_{\Gamma^{\prime}}e^{\lambda^{\prime}s}(\lambda-\lambda^{\prime})^{-1}\,\mathrm{d}\lambda^{\prime}\,\mathrm{d}\lambda\Bigg]\\
&=\dfrac{1}{2\pi i}\int_\Gamma e^{\lambda(t+s)}(\lambda I-A)^{-1}\,\mathrm{d}\lambda=T(t+s).
\end{aligned}\]Here we have used the identity $(\lambda I-A)^{-1}(\lambda^{\prime}I-A)^{-1}=(\lambda-\lambda^{\prime})^{-1}[(\lambda^{\prime}I-A)^{-1}-(\lambda I-A)^{-1}]$ and the fact that\[\begin{aligned}
\int_\Gamma e^{\lambda t}(\lambda-\lambda^{\prime})^{-1}\,\mathrm{d}\lambda&=0,\\
\int_{\Gamma^{\prime}}e^{\lambda^{\prime}s}(\lambda-\lambda^{\prime})^{-1}\,\mathrm{d}\lambda^{\prime}&=-2\pi i e^{\lambda s}.
\end{aligned}\]

By shifting the contour to the right if necessary, we may assume that $\Gamma$ in \eqref{eq:analytic-semigroup-contour-integral-representation} is entirely in the set $\{|\arg\lambda|<\phi\}$. For sufficiently small $t$, we may then make the substitution $\lambda^{\prime}=\lambda|t|$ and then deform the contour back to $\Gamma$. In this fashion, we obtain
\begin{equation}\label{eq:analytic-semigroup-contour-rescaled-at-small-time}
T(t)=\dfrac{1}{2\pi i|t|}\int_\Gamma e^{\lambda^{\prime}t/|t|}\left(\dfrac{\lambda^{\prime}}{|t|}I-A\right)^{-1}\,\mathrm{d}\lambda^{\prime}.
\end{equation}
From this, we find
\begin{equation}\label{eq:analytic-semigroup-small-time-contour-norm-bound}
\|T(t)\|\leq\dfrac{C}{2\pi}\int_\Gamma|e^{\lambda^{\prime}t/|t|}||\lambda^{\prime}|^{-1}\,\mathrm{d}\lambda^{\prime}.
\end{equation}
Clearly, this expression remains bounded as $t\to0$. Moreover, for $u\in\mathcal{D}(A)$, we find, if we choose $\Gamma$ lying to the right of the origin,\[\begin{aligned}
T(t)u-u&=\dfrac{1}{2\pi i}\int_\Gamma e^{\lambda t}((\lambda I-A)^{-1}-\lambda^{-1})u\,\mathrm{d}\lambda\\
&=-\dfrac{1}{2\pi i}\int_\Gamma e^{\lambda t}\lambda^{-1}(\lambda I-A)^{-1}\,\mathrm{d}\lambda\ Au.
\end{aligned}\]As $t\to0$, the integral converges to
\begin{equation}\label{eq:analytic-semigroup-origin-contour-zero-integral}
-\dfrac{1}{2\pi i}\int_\Gamma\lambda^{-1}(\lambda I-A)^{-1}\,\mathrm{d}\lambda=0,
\end{equation}
as can be seen by closing the contour by a circle on the right. This shows that $T(t)$ converges strongly to the identity as $t\to0$.

\begingroup
\renewenvironment{Proof}{\par}{\par\nobreak\hfill$\blacksquare$\par\vspace{0.1cm}}
\begin{Proof}
For $t\neq0$, we find from \eqref{eq:analytic-semigroup-contour-integral-representation}:\[\begin{aligned}
\dfrac{\mathrm{d}}{\mathrm{d}t}T(t)&=\dfrac{1}{2\pi i}\int_\Gamma\lambda e^{\lambda t}(\lambda I-A)^{-1}\,\mathrm{d}\lambda\\
&=\dfrac{1}{2\pi i}\int_\Gamma e^{\lambda t}A(\lambda I-A)^{-1}\,\mathrm{d}\lambda+\dfrac{1}{2\pi i}\int_\Gamma e^{\lambda t}I\,\mathrm{d}\lambda\\
&=AT(t).
\end{aligned}\]Here the closedness of $A$ is used to justify taking it outside the integral. Hence we have\[\begin{aligned}
\dfrac{T(h)-I}{h}u&=\lim_{\epsilon\to0}\dfrac{T(h)-T(\epsilon)}{h}u\\
&=\lim_{\epsilon\to0}\dfrac{1}{h}\int_\epsilon^h AT(t)u\,\mathrm{d}t=\dfrac{1}{h}\int_0^h T(t)Au
\end{aligned}\]for $u\in\mathcal{D}(A)$. As $h\to0$, we obtain $Au$ in the limit; hence the infinitesimal generator of the semigroup is an extension of $A$. Since $A$ has a resolvent in the half-plane $\operatorname{Re}\lambda>\omega$, this extension cannot be proper.
\end{Proof}
\endgroup

In the course of the proof above, we proved that $\exp(At)$ is differentiable for $t>0$ and that $\frac{\mathrm{d}}{\mathrm{d}t}\exp(At)=A\exp(At)$. Using in essence the same argument, we can show that $\exp(At)$ is actually infinitely often differentiable and $\frac{\mathrm{d}^n}{\mathrm{d}t^n}\exp(At)=A^n\exp(At)$. In particular, this implies that the range of $\exp(At)$ is contained in $\mathcal{D}(A^n)$ for every $n$. In PDE applications, this usually translates into a smoothing property. Even if the initial conditions have singularities, the solution is smooth (of class $C^\infty$) for positive $t$. For future use, we note the following bound on the norm of $A\exp(At)$.

\begin{Lemma}
\label{lem:analytic-semigroup-time-derivative-bound}
Let $A$ be the infinitesimal generator of an analytic semigroup. Then there are constants $C$ and $\omega$ such that, for every $t>0$, we have
\begin{equation}\label{eq:analytic-semigroup-generator-time-derivative-bound}
\|A\exp(At)\|\leq C\exp(\omega t)/t.
\end{equation}

\end{Lemma}

The proof is based on a calculation similar to that leading to \eqref{eq:analytic-semigroup-small-time-contour-norm-bound}; we leave the details as an exercise (Problem~\ref{ex:analytic-semigroup-derivative-bound-proof}).

In ODE courses we learn that the zero solution of the equation $\dot u=Au$ is stable if all eigenvalues of $A$ have negative real parts. Indeed, this is the most widely applied practical criterion for evaluating stability. Naturally, one would like to apply the same type of criterion for physical systems described by PDEs. This raises the following question: If $A$ is the generator of a semigroup and the spectrum of $A$ is entirely in the left half-plane at a positive distance from the imaginary axis, does it follow that $\|\exp(At)\|$ is bounded as $t\to\infty$? Or equivalently, is there a connection between an upper bound on the real part of the spectrum of $A$ and the type of the semigroup? Unfortunately, in the general context of $C_0$-semigroups, the answer is negative. Although the counterexamples are not the kind that would arise ``naturally\textquotesingle\textquotesingle\ in physical applications, no practical conditions are known that would lead to a positive result. The state of current knowledge is clearly unsatisfactory. For analytic semigroups, however, the problem is quite easy. We have the following result.

\begin{Theorem}
\label{thm:analytic-semigroup-spectrum-controls-growth}
Let $A$ be the infinitesimal generator of an analytic semigroup and assume that the spectrum of $A$ is entirely to the left of the line $\operatorname{Re}\lambda=\omega$. Then there exists a constant $M$ such that $\|\exp(At)\|\leq M\exp(\omega t)$.
\end{Theorem}

The idea of the proof is to shift the contour in \eqref{eq:analytic-semigroup-contour-integral-representation} so that it lies entirely to the left of $\operatorname{Re}\lambda=\omega$. We leave the details of the argument as an exercise (Problem~\ref{ex:analytic-semigroup-spectral-growth-proof}).

\subsubsection{Fractional Powers}\label{subsec:analytic-generator-fractional-powers}

Throughout this subsection, we assume that $A$ is the infinitesimal generator of an analytic semigroup and that the spectrum of $A$ lies entirely in the (open) left half-plane. In this case, we shall define fractional powers of $-A$. If the spectrum of $A$ lies in the half-plane $\operatorname{Re}\lambda<\omega$, then we can of course apply the same considerations to define fractional powers of $\omega I-A$. If the spectrum is in the open left half-plane, we can choose a $\delta>0$ such that the spectrum of $A$ is in the half-plane $\operatorname{Re}\lambda<-\delta$, and from the previous subsection we have bounds of the form $\|\exp(At)\|\leq M\exp(-\delta t)$, $\|A\exp(At)\|\leq M_1\exp(-\delta t)/t$. Moreover, by noting that $A^n\exp(At)=(A\exp(A\frac{t}{n}))^n$, we find $\|A^n\exp(At)\|\leq M_n\exp(-\delta t)/t^n$.

We define
\begin{equation}\label{eq:analytic-generator-negative-fractional-power-contour}
(-A)^{-\alpha}=-\dfrac{1}{2\pi i}\int_\Gamma\lambda^{-\alpha}(\lambda I+A)^{-1}\,\mathrm{d}\lambda,
\end{equation}
where $\Gamma$ is a curve from $\exp(-i\theta)\infty$ to $\exp(i\theta)\infty$ such that the spectrum of $-A$ lies to the right and the origin lies to the left of $\Gamma$. Here $\pi/2-\delta<\theta<\pi$, where $\delta$ is as in Theorem~\ref{thm:analytic-generator-sectorial-resolvent} and $\lambda^{-\alpha}$ denotes the branch of the function which takes positive values on the positive real axis. It follows from \eqref{eq:analytic-semigroup-half-plane-resolvent-estimate} that the integral in \eqref{eq:analytic-generator-negative-fractional-power-contour} is absolutely convergent for any $\alpha>0$. If $\alpha$ is an integer, we can deform the contour in \eqref{eq:analytic-generator-negative-fractional-power-contour} to a circle around the origin and use residues to evaluate the integral; we find that indeed $(-A)^{-n}=(-A^{-1})^n$; see Problem~\ref{ex:fractional-power-integer-inverse-identity}. If $0<\alpha<1$, we can deform $\Gamma$ into the upper and lower sides of the negative real axis. This leads to the expression
\begin{equation}\label{eq:analytic-generator-negative-fractional-power-real-resolvent}
(-A)^{-\alpha}=\dfrac{\sin(\pi\alpha)}{\pi}\int_0^\infty\lambda^{-\alpha}(\lambda I-A)^{-1}\,\mathrm{d}\lambda,\ 0<\alpha<1.
\end{equation}


Recall from the proof of the Hille-Yosida theorem that
\begin{equation}\label{eq:fractional-power-resolvent-laplace-formula}
(\lambda I-A)^{-1}=\int_0^\infty e^{-\lambda t}e^{At}\,\mathrm{d}t.
\end{equation}
We insert this expression into \eqref{eq:analytic-generator-negative-fractional-power-real-resolvent} and then exchange the order of integrations. This yields\[\begin{aligned}
(-A)^{-\alpha}&=\dfrac{\sin(\pi\alpha)}{\pi}\int_0^\infty\exp(At)\int_0^\infty\lambda^{-\alpha}e^{-\lambda t}\,\mathrm{d}\lambda\,\mathrm{d}t\\
&=\dfrac{\sin(\pi\alpha)}{\pi}\left(\int_0^\infty u^{-\alpha}e^{-u}\,\mathrm{d}u\right)\left(\int_0^\infty t^{\alpha-1}e^{At}\,\mathrm{d}t\right).
\end{aligned}\]We have
\begin{equation}\label{eq:fractional-power-gamma-reflection-integral}
\int_0^\infty u^{-\alpha}e^{-u}\,\mathrm{d}u=\dfrac{\pi}{\sin(\pi\alpha)\Gamma(\alpha)},
\end{equation}
and hence
\begin{equation}\label{eq:analytic-generator-fractional-power-semigroup-integral}
(-A)^{-\alpha}=\dfrac{1}{\Gamma(\alpha)}\int_0^\infty t^{\alpha-1}e^{At}\,\mathrm{d}t.
\end{equation}
The argument we gave here applies only if $0<\alpha<1$. However, both \eqref{eq:analytic-generator-negative-fractional-power-contour} and \eqref{eq:analytic-generator-fractional-power-semigroup-integral} are defined for any $\alpha>0$ and analytic in $\alpha$, hence by the uniqueness of analytic continuation they agree for every $\alpha>0$. Using \eqref{eq:analytic-generator-fractional-power-semigroup-integral} and the bound $\|\exp(At)\|\leq M\exp(-\delta t)$, one easily establishes the following result.

\begin{Lemma}
\label{lem:negative-fractional-power-uniform-bound}
There is a constant $C$ such that $\|(-A)^{-\alpha}\|\leq C$ for $0<\alpha\leq1$.
\end{Lemma}

\begin{Theorem}
\label{thm:negative-fractional-powers-semigroup}
The following hold.\begin{enumerate}\item $(-A)^{-\alpha}(-A)^{-\beta}=(-A)^{-(\alpha+\beta)}$.\item $\lim_{\alpha\to0}(-A)^{-\alpha}=I$ in the strong operator topology.\end{enumerate}
\end{Theorem}

With the obvious convention $(-A)^0=I$, this theorem asserts that the operators $(-A)^{-\alpha}$, $\alpha\geq0$, form a $C_0$-semigroup. It is natural to call the infinitesimal generator $-\log(-A)$, but we shall not pursue this point further.

\begin{Proof}
We have\[\begin{aligned}
&(-A)^{-\alpha}(-A)^{-\beta}\\
&\quad=\dfrac{1}{\Gamma(\alpha)\Gamma(\beta)}\int_0^\infty\int_0^\infty t^{\alpha-1}s^{\beta-1}e^{At}e^{As}\,\mathrm{d}t\,\mathrm{d}s\\
&\quad=\dfrac{1}{\Gamma(\alpha)\Gamma(\beta)}\int_0^\infty t^{\alpha-1}\int_t^\infty(u-t)^{\beta-1}e^{Au}\,\mathrm{d}u\,\mathrm{d}t\\
&\quad=\dfrac{1}{\Gamma(\alpha)\Gamma(\beta)}\int_0^\infty\int_0^u t^{\alpha-1}(u-t)^{\beta-1}\,\mathrm{d}t\ e^{Au}\,\mathrm{d}u\\
&\quad=\dfrac{1}{\Gamma(\alpha)\Gamma(\beta)}\int_0^1 v^{\alpha-1}(1-v)^{\beta-1}\,\mathrm{d}v\int_0^\infty u^{\alpha+\beta-1}e^{Au}\,\mathrm{d}u\\
&\quad=\dfrac{1}{\Gamma(\alpha+\beta)}\int_0^\infty u^{\alpha+\beta-1}e^{Au}\,\mathrm{d}u\\
&\quad=(-A)^{-\alpha-\beta}.
\end{aligned}\]Since we already know that $\|(-A)^{-\alpha}\|$ is bounded as $\alpha\to0$, it suffices to show that $(-A)^{-\alpha}u\to u$ for $u$ in a dense subset of $X$. Choose $u\in\mathcal{D}(A)$, then $u=-A^{-1}y$ for some $y\in X$. We then have
\begin{equation}\label{eq:fractional-power-semigroup-strong-continuity-reduction}
(-A)^{-\alpha}u-u=(-A)^{-1-\alpha}y-(-A)^{-1}y,
\end{equation}
and it is clear from either \eqref{eq:analytic-generator-negative-fractional-power-contour} or \eqref{eq:analytic-generator-fractional-power-semigroup-integral} that for $\alpha>0$, $(-A)^{-\alpha}$ is actually continuous (indeed analytic) in the uniform operator topology.
\end{Proof}

Since $(-A)^{-n}$ is one-to-one for $n\in\mathbb{N}$ and $(-A)^{-n}=(-A)^{-n+\alpha}(-A)^{-\alpha}$ for $n>\alpha$, it follows that $(-A)^{-\alpha}$ is one-to-one. Hence it has an inverse, which we naturally denote by $(-A)^\alpha$. It is clear that $(-A)^\alpha$ is closed with domain $\mathcal{D}((-A)^\alpha)=\mathcal{R}((-A)^{-\alpha})$; since $\mathcal{R}((-A)^{-n})=\mathcal{D}(A^n)\subset\mathcal{R}((-A)^{-\alpha})$ for $n>\alpha$, it follows that the domain of $(-A)^\alpha$ is dense. Moreover, it is easy to check that $(-A)^{\alpha+\beta}u=(-A)^\alpha(-A)^\beta u$ for any $\alpha,\beta\in\mathbb{R}$ and any $u\in\mathcal{D}((-A)^\gamma)$, where $\gamma=\max(\alpha,\beta,\alpha+\beta)$.\par We conclude this subsection with a result relating $(-A)^\alpha$ to the semigroup.

\begin{Lemma}
\label{lem:fractional-power-semigroup-smoothing-holder}
Let $\alpha>0$. For every $u\in\mathcal{D}((-A)^\alpha)$, we have\[\exp(At)(-A)^\alpha u=(-A)^\alpha\exp(At)u.\]Moreover, the operator $(-A)^\alpha\exp(At)$ is bounded, with a bound of the form
\begin{equation}\label{eq:analytic-semigroup-fractional-power-smoothing-bound}
\|(-A)^\alpha\exp(At)\|\leq M_\alpha t^{-\alpha}e^{-\delta t}.
\end{equation}
If $0<\alpha\leq1$ and $u\in\mathcal{D}((-A)^\alpha)$, we have a bound of the form
\begin{equation}\label{eq:analytic-semigroup-fractional-domain-holder-orbit-bound}
\|\exp(At)u-u\|\leq C_\alpha t^\alpha\|(-A)^\alpha u\|.
\end{equation}

\end{Lemma}

\begin{Proof}
If $u\in\mathcal{D}((-A)^\alpha)$, then $u=(-A)^{-\alpha}v$ for some $v$, and we find\[\begin{aligned}
\exp(At)u&=\exp(At)(-A)^{-\alpha}v=\dfrac{1}{\Gamma(\alpha)}\int_0^\infty s^{\alpha-1}\exp(As)\exp(At)v\,\mathrm{d}s\\
&=(-A)^{-\alpha}\exp(At)v=(-A)^{-\alpha}\exp(At)(-A)^\alpha u.
\end{aligned}\]The first claim of the lemma follows by applying $(-A)^\alpha$ to both sides.\par Let $n-1<\alpha<n$, then\[\begin{aligned}
\|(-A)^\alpha\exp(At)\|&=\|(-A)^{\alpha-n}A^n\exp(At)\|\\
&\leq\dfrac{1}{\Gamma(n-\alpha)}\int_0^\infty s^{n-\alpha-1}\|A^n\exp(A(t+s))\|\,\mathrm{d}s\\
&\leq\dfrac{M_n}{\Gamma(n-\alpha)}\int_0^\infty s^{n-\alpha-1}(t+s)^{-n}e^{-\delta(t+s)}\,\mathrm{d}s\\
&\leq\dfrac{M_ne^{-\delta t}}{\Gamma(n-\alpha)t^\alpha}\int_0^\infty r^{n-\alpha-1}(1+r)^{-n}\,\mathrm{d}r.
\end{aligned}\]Finally, we have\[\begin{aligned}
\|e^{At}u-u\|&=\left\|\int_0^t Ae^{As}u\,\mathrm{d}s\right\|=\left\|\int_0^t(-A)^{1-\alpha}e^{As}(-A)^\alpha u\,\mathrm{d}s\right\|\\
&\leq C\int_0^t s^{\alpha-1}\|(-A)^\alpha u\|\,\mathrm{d}s\leq C_\alpha t^\alpha\|(-A)^\alpha u\|.
\end{aligned}\]This completes the proof.
\end{Proof}

\subsubsection{Perturbations of Analytic Semigroups}\label{subsec:analytic-semigroup-perturbations}

Suppose $A$ is the infinitesimal generator of a semigroup. In what sense must $B$ be ``small\textquotesingle\textquotesingle\ to be sure that $A+B$ also generates a semigroup? In the context of $C_0$-semigroups, we saw in Problem~\ref{ex:quasicontraction-bounded-operator-perturbation} that it suffices if $B$ is bounded. Indeed, unless conditions other than smallness are imposed, e.g., that both $A$ and $B$ are dissipative, then this is basically the best we can do. In particular, it does not help if $B$ is ``of lower order.\textquotesingle\textquotesingle\ The reader may verify that although the operator $Au=u_{xxx}$ generates a $C_0$-semigroup in $L^2(\mathbb{R})$, the operator $(A+B)u=u_{xxx}-u_{xx}$ does not. For analytic semigroups, however, the situation is much better. The main result is the following theorem.

\begin{Theorem}
\label{thm:analytic-generator-small-relative-perturbation}
Let $A$ be the infinitesimal generator of an analytic semigroup. Then there exists a positive number $\delta$ such that, if $B$ is any operator satisfying\begin{enumerate}\item $B$ is closed and $\mathcal{D}(B)\supseteq\mathcal{D}(A)$,\item $\|Bu\|\leq a\|Au\|+b\|u\|$ for $u\in\mathcal{D}(A)$, where $a<\delta$,\end{enumerate}then $A+B$ is also the infinitesimal generator of an analytic semigroup.
\end{Theorem}

\begin{Proof}
Since $A$ generates an analytic semigroup, there exists $\omega\in\mathbb{R}$ and $M>0$ such that $\|R_\lambda(A)\|\leq M/|\lambda-\omega|$ for $\operatorname{Re}\lambda>\omega$. The operator $BR_\lambda(A)$ is bounded, and we find\[\begin{aligned}
\|BR_\lambda(A)u\|&\leq a\|AR_\lambda(A)u\|+b\|R_\lambda(A)u\|\\
&\leq a\left(1+\dfrac{M|\lambda|}{|\lambda-\omega|}\right)\|u\|+\dfrac{bM}{|\lambda-\omega|}\|u\|.
\end{aligned}\]For any $\epsilon>0$, we can find $\omega^{\prime}$ such that $\|BR_\lambda(A)\|\leq a(1+M+\epsilon)$ for $\operatorname{Re}\lambda>\omega^{\prime}$. If moreover $a<(1+M)^{-1}$, then we can choose $\epsilon$ such that $\|BR_\lambda(A)\|<1$ for $\operatorname{Re}\lambda>\omega^{\prime}$. The rest follows from the identity
\begin{equation}\label{eq:analytic-generator-perturbation-resolvent-identity}
R_\lambda(A+B)=R_\lambda(A)(I+BR_\lambda(A))^{-1},
\end{equation}
from which we find $\|R_\lambda(A+B)\|\leq M^{\prime}/|\lambda-\omega^{\prime}|$ for $\operatorname{Re}\lambda>\omega^{\prime}$.
\end{Proof}

In applications, $B$ is often ``of lower order\textquotesingle\textquotesingle\ than $A$, and $a$ in the last theorem can be taken arbitrarily small. The abstract form of the notion of ``lower order\textquotesingle\textquotesingle\ can be phrased in terms of fractional powers. We have the following lemma.

\begin{Lemma}
\label{lem:analytic-fractional-lower-order-relative-bound}
Let $A$ be the infinitesimal generator of an analytic semigroup and assume that $B$ is closed and $\mathcal{D}(B)\supseteq\mathcal{D}((\omega I-A)^\alpha)$ for some $\alpha\in(0,1)$. Then there is a constant $C$ such that
\begin{equation}\label{eq:analytic-generator-fractional-lower-order-relative-bound}
\|Bu\|\leq C(\rho^\alpha\|u\|+\rho^{\alpha-1}\|(A-\omega I)u\|)
\end{equation}
for every $u\in\mathcal{D}(A)$ and every $\rho>0$.
\end{Lemma}

By choosing $\rho$ sufficiently large and applying the last theorem, we conclude that $A+B$ generates an analytic semigroup.

\begin{Proof}
Without loss of generality, we may assume $\omega=0$. If $\mathcal{D}(B)\supseteq\mathcal{D}((-A)^\alpha)$, then $B(-A)^{-\alpha}$ is bounded; i.e., there is a constant $C$ such that $\|Bu\|\leq C\|(-A)^\alpha u\|$. Hence it suffices to show \eqref{eq:analytic-generator-fractional-lower-order-relative-bound} for $B=(-A)^\alpha$. We have, for $u\in\mathcal{D}(A)$,\[\begin{aligned}
\|(-A)^\alpha u\|&\leq\left|\dfrac{\sin(\pi\alpha)}{\pi}\right|\int_0^\rho\lambda^{\alpha-1}\|AR_\lambda(A)\|\|u\|\,\mathrm{d}\lambda\\
&\quad+\left|\dfrac{\sin(\pi\alpha)}{\pi}\right|\int_\rho^\infty\lambda^{\alpha-1}\|R_\lambda(A)\|\|Au\|\,\mathrm{d}\lambda.
\end{aligned}\]We now use the fact that $\|R_\lambda(A)\|\leq M/\lambda$ and $\|AR_\lambda(A)\|\leq1+M$ to complete the proof of the lemma.
\end{Proof}

In applications, it is often difficult to precisely characterize the domains of fractional powers. Instead of checking that $\mathcal{D}(B)\supseteq\mathcal{D}((\omega I-A)^\alpha)$, one usually checks \eqref{eq:analytic-generator-fractional-lower-order-relative-bound} directly. In this context, the following result is of interest.

\begin{Lemma}
\label{lem:analytic-relative-estimate-implies-fractional-domain}
Let $A$ be the generator of an analytic semigroup and let $B$ be a closed linear operator such that $\mathcal{D}(B)\supseteq\mathcal{D}(A)$ and, for some $\gamma\in(0,1)$ and every $\rho\geq\rho_0>0$, we have
\begin{equation}\label{eq:analytic-generator-relative-estimate-fractional-domain-hypothesis}
\|Bu\|\leq C(\rho^\gamma\|u\|+\rho^{\gamma-1}\|Au\|)
\end{equation}
for every $u\in\mathcal{D}(A)$. Then $\mathcal{D}(B)\supseteq\mathcal{D}((\omega-A)^\alpha)$ for every $\alpha>\gamma$.
\end{Lemma}

\begin{Proof}
Again we assume without loss of generality that $\omega=0$. Let $u\in\mathcal{D}((-A)^{1-\alpha})$ so that $(-A)^{-\alpha}u\in\mathcal{D}(A)$. We have
\begin{equation}\label{eq:analytic-generator-fractional-domain-converse-integral}
B(-A)^{-\alpha}u=\dfrac{1}{\Gamma(\alpha)}\int_0^\infty t^{\alpha-1}B\exp(At)\,\mathrm{d}t,
\end{equation}
provided that the integral is convergent. We split the integral as
\begin{equation}\label{eq:analytic-generator-fractional-domain-converse-integral-split}
\int_0^\delta t^{\alpha-1}B\exp(At)\,\mathrm{d}t+\int_\delta^\infty t^{\alpha-1}B\exp(At)\,\mathrm{d}t.
\end{equation}
We set $\delta=1/\rho_0$ and use \eqref{eq:analytic-generator-relative-estimate-fractional-domain-hypothesis} with $\rho=\rho_0$ in the second integral and $\rho=1/t$ in the first integral. The result is that $B(-A)^{-\alpha}$ is bounded for $\alpha>\gamma$, which implies the lemma.
\end{Proof}

We now present an application to parabolic PDEs. Let $\Omega$ be a bounded domain in $\mathbb{R}^m$ with smooth boundary, let $a_{ij}(\bm{x})$ be of class $C^1(\overline\Omega)$ such that the matrix $a_{ij}$ is symmetric and strictly positive definite and let $b_i(\bm{x})$, $c(\bm{x})$ be of class $C^0(\overline\Omega)$. In $L^2(\Omega)$, we consider the operator
\begin{equation}\label{eq:analytic-semigroup-elliptic-dirichlet-generator}
Au=\dfrac{\partial}{\partial x_i}\left(a_{ij}(\bm{x})\dfrac{\partial u}{\partial x_j}\right)+b_i(\bm{x})\dfrac{\partial u}{\partial x_i}+c(\bm{x})u
\end{equation}
with domain $H^2(\Omega)\cap H_0^1(\Omega)$. We claim

\begin{Theorem}
\label{thm:elliptic-dirichlet-analytic-generation}
$A$ generates an analytic semigroup.
\end{Theorem}

\begin{Proof}
Let
\begin{equation}\label{eq:analytic-semigroup-elliptic-principal-self-adjoint-part}
A_0u=\dfrac{\partial}{\partial x_i}\left(a_{ij}(\bm{x})\dfrac{\partial u}{\partial x_j}\right).
\end{equation}
Then $A_0$ is self-adjoint with negative spectrum; hence it clearly generates an analytic semigroup. Moreover, we find\[\begin{aligned}
\|(A-A_0)u\|_2&\leq C\|u\|_{1,2}\leq C\|u\|_2^{1/2}\|u\|_{2,2}^{1/2}\\
&\leq C\|u\|_2^{1/2}\|A_0u\|_2^{1/2}\leq C(\rho^{1/2}\|u\|_2+\rho^{-1/2}\|A_0u\|_2).
\end{aligned}\]Hence $\mathcal{D}(A-A_0)$ contains $\mathcal{D}((-A_0)^\alpha)$ for any $\alpha>1/2$.
\end{Proof}

\begin{Remark}
\label{rem:elliptic-square-root-dirichlet-domain}
The intelligent reader may suspect that $\mathcal{D}((-A_0)^{1/2})$ is actually $H_0^1(\Omega)$. Indeed, this suspicion is well founded. A proof, however, would be significantly more involved than the discussion given above.
\end{Remark}

\subsubsection{Regularity of Mild Solutions}\label{subsec:analytic-mild-solution-regularity}

We now turn our attention to the inhomogeneous initial-value problem
\begin{equation}\label{eq:analytic-semigroup-inhomogeneous-mild-initial-value-problem}
\dot u(t)=Au(t)+f(t),\ u(0)=u_0.
\end{equation}
The mild solution is given by
\begin{equation}\label{eq:analytic-semigroup-mild-solution-representation}
u(t)=e^{At}u_0+\int_0^t e^{A(t-s)}f(s)\,\mathrm{d}s.
\end{equation}
If $A$ generates an analytic semigroup, we already know that the term $e^{At}u_0$ is analytic in $t$ for $t>0$; moreover, $e^{At}u_0$ is in $\mathcal{D}(A^n)$ for every $n$. Moreover, we know that $\|A^ne^{At}u_0\|\leq C\|u_0\|/t^n$ as $t\to0$. We can hence focus attention on the term
\begin{equation}\label{eq:analytic-semigroup-mild-inhomogeneous-convolution}
v(t):=\int_0^t e^{A(t-s)}f(s)\,\mathrm{d}s.
\end{equation}
We need the following definition:

\begin{Definition}
\label{def:banach-valued-holder-continuity}
We say that $f\in C^\theta([0,T];X)$, $0<\theta<1$, if there is a constant $L$ such that
\begin{equation}\label{eq:banach-valued-holder-continuity-definition}
\|f(t)-f(s)\|\leq L|t-s|^\theta\quad\forall s,t\in[0,T].
\end{equation}

\end{Definition}

\begin{Lemma}
\label{lem:analytic-holder-data-difference-convolution}
Let $A$ be the infinitesimal generator of an analytic semigroup in $X$ and assume that $f\in C^\theta([0,T];X)$ for some $\theta\in(0,1)$. Let
\begin{equation}\label{eq:analytic-semigroup-holder-data-difference-convolution}
w(t)=\int_0^t e^{A(t-s)}(f(s)-f(t))\,\mathrm{d}s.
\end{equation}
Then $w(t)\in\mathcal{D}(A)$ for every $t\in[0,T]$ and $Aw\in C^\theta([0,T];X)$.
\end{Lemma}

\begin{Proof}
Let us assume that $\|\exp(At)\|\leq M$ and $\|A\exp(At)\|\leq C/t$ for $t\in(0,T]$. The fact that $w(t)\in\mathcal{D}(A)$ follows from the estimate
\begin{equation}\label{eq:analytic-holder-data-difference-in-generator-domain-bound}
\left\|\int_0^t Ae^{A(t-s)}(f(s)-f(t))\,\mathrm{d}s\right\|\leq\int_0^t\dfrac{C}{t-s}L(t-s)^\theta\,\mathrm{d}s=\dfrac{CLt^\theta}{\theta}.
\end{equation}
It remains to prove the Hölder continuity. We first note that
\begin{equation}\label{eq:analytic-generator-semigroup-time-difference-bound}
\begin{aligned}\|A\exp(At)-A\exp(As)\|&=\left\|\int_s^t A^2\exp(A\tau)\,\mathrm{d}\tau\right\|\\&\leq\int_s^t\|A^2\exp(A\tau)\|\,\mathrm{d}\tau\\&\leq4C\int_s^t\tau^{-2}\,\mathrm{d}\tau=4Ct^{-1}s^{-1}(t-s).\end{aligned}
\end{equation}
We next write\[\begin{aligned}
&Aw(t+h)-Aw(t)\\
&\quad=A\int_0^t[\exp(A(t+h-s))-\exp(A(t-s))][f(s)-f(t)]\,\mathrm{d}s\\
&\qquad+A\int_0^t\exp(A(t+h-s))(f(t)-f(t+h))\,\mathrm{d}s\\
&\qquad+A\int_t^{t+h}\exp(A(t+h-s))(f(s)-f(t+h))\,\mathrm{d}s\\
&\quad=:I_1+I_2+I_3.
\end{aligned}\]We now use \eqref{eq:analytic-generator-semigroup-time-difference-bound} to obtain
\begin{equation}\label{eq:analytic-holder-difference-convolution-first-term-bound}
\|I_1\|\leq\int_0^t4C(t+h-s)^{-1}(t-s)^{-1}hL(t-s)^\theta\,\mathrm{d}s\leq C_1h^\theta.
\end{equation}
(Use the substitution $s=t-h\tau$.) $I_2$ can be rewritten as $(\exp(A(t+h))-\exp(At))(f(t)-f(t+h))$, and the Hölder estimate simply follows from that for the second factor. For the last term, we have
\begin{equation}\label{eq:analytic-holder-difference-convolution-third-term-bound}
\|I_3\|\leq\int_t^{t+h}\dfrac{C}{t+h-s}L(t+h-s)^\theta\,\mathrm{d}s\leq C_2h^\theta.
\end{equation}

\end{Proof}

The following regularity result says that, except for a neighborhood of $t=0$, $\dot u$ and $Au$ are as smooth as $f$ is. Note that there is no comparable result for $C_0$-semigroups.

\begin{Theorem}
\label{thm:analytic-mild-holder-regularity}
Let $A$ be the infinitesimal generator of an analytic semigroup and let $u$ be the solution of \eqref{eq:analytic-semigroup-inhomogeneous-mild-initial-value-problem} as given by \eqref{eq:analytic-semigroup-mild-solution-representation}. Moreover, let $f\in C^\theta([0,T];X)$. Then:\begin{enumerate}\item\label{concl:analytic-mild-regularity-away-from-zero} For every $\delta>0$, $Au$ and $\dot u$ are in $C^\theta([\delta,T];X)$.\item\label{concl:analytic-mild-regularity-continuity-at-zero} If $u_0\in\mathcal{D}(A)$, then $Au$ and $\dot u$ are in $C([0,T];X)$.\item\label{concl:analytic-mild-regularity-zero-compatible-holder} If $u_0=0$ and $f(0)=0$, then $Au$ and $\dot u$ are in $C^\theta([0,T];X)$.\end{enumerate}
\end{Theorem}

\begin{Proof}
Let $v(t)$ be as given by \eqref{eq:analytic-semigroup-mild-inhomogeneous-convolution}. Recall from the proof of Theorem~\ref{thm:semigroup-mild-solution-classical-regularity} that if $Av(t)$ exists, then $\dot v$ also exists and $\dot v=Av+f$.\par Hence, it suffices to consider $Au$ in verifying the theorem. For this purpose, we write
\begin{equation}\label{eq:analytic-mild-regularity-three-part-solution-splitting}
u(t)=e^{At}u_0+\int_0^t e^{A(t-s)}(f(s)-f(t))\,\mathrm{d}s+\int_0^t e^{A(t-s)}f(t)\,\mathrm{d}s.
\end{equation}
In view of the previous lemma, it suffices to consider the last term. We have
\begin{equation}\label{eq:analytic-mild-regularity-frozen-data-generator-integral}
A\int_0^t e^{A(t-s)}f(t)\,\mathrm{d}s=(e^{At}-I)f(t).
\end{equation}
Since $f$ was assumed in $C^\theta([0,T];X)$, we need only consider $\exp(At)f(t)$. We have, for $t\geq\delta$ and $h>0$,
\begin{equation}\label{eq:analytic-mild-regularity-holder-estimate-away-from-zero}
\begin{aligned}[t]&\|\exp(A(t+h))f(t+h)-\exp(At)f(t)\|\\&\quad\leq\|\exp(A(t+h))\|\ \|f(t+h)-f(t)\|\\&\qquad+\|\exp(A(t+h))-\exp(At)\|\ \|f(t)\|\\&\quad\leq C_1h^\theta+C_2\dfrac{h}{\delta}.\end{aligned}
\end{equation}
This implies part~\ref{concl:analytic-mild-regularity-away-from-zero}. Next we note that
\begin{equation}\label{eq:analytic-mild-regularity-initial-strong-continuity-bound}
\begin{aligned}\|\exp(At)f(t)-f(0)\|&\leq\|\exp(At)f(0)-f(0)\|\\&\quad+\|\exp(At)\|\|f(t)-f(0)\|;\end{aligned}
\end{equation}
the strong continuity of the semigroup now implies part~\ref{concl:analytic-mild-regularity-continuity-at-zero}. To show part~\ref{concl:analytic-mild-regularity-zero-compatible-holder}, we first proceed as in \eqref{eq:analytic-mild-regularity-holder-estimate-away-from-zero}, but then we estimate\[\begin{aligned}
\|(e^{A(t+h)}-e^{At})f(t)\|&=\left\|\int_t^{t+h}Ae^{A\tau}f(t)\,\mathrm{d}\tau\right\|\\
&\leq\int_t^{t+h}\|Ae^{A\tau}(f(t)-f(0))\|\,\mathrm{d}\tau\\
&\leq C\int_t^{t+h}\tau^{-1}t^\theta\,\mathrm{d}\tau\\
&\leq C\int_t^{t+h}\tau^{\theta-1}\,\mathrm{d}\tau\leq Ch^\theta.
\end{aligned}\]This completes the proof
\end{Proof}

\Needspace{5\baselineskip}\textbf{Problems}

\begin{enumerate}[label=\thesection.\arabic*.,ref=\thesection.\arabic*,start=20,series=book-problems]

\item\label{ex:analytic-semigroup-derivative-bound-proof} Prove Lemma~\ref{lem:analytic-semigroup-time-derivative-bound}.

\item\label{ex:analytic-semigroup-spectral-growth-proof} Prove Theorem~\ref{thm:analytic-semigroup-spectrum-controls-growth}.

\item\label{ex:fractional-power-integer-inverse-identity} Verify that for $n\in\mathbb{N}$, we have $(-A)^{-n}=(-A^{-1})^n$, where $(-A)^{-n}$ is defined by \eqref{eq:analytic-generator-negative-fractional-power-contour}.

\item\label{ex:analytic-mild-fractional-compatible-initial-regularity} Prove part~\ref{concl:analytic-mild-regularity-zero-compatible-holder} of Theorem~\ref{thm:analytic-mild-holder-regularity}, assuming that $u_0$ and $f(0)$ lie in the domain of appropriate fractional powers of $\omega-A$.

\item\label{ex:analytic-semigroup-off-diagonal-product-space-generator} Let $A$ be the infinitesimal generator of an analytic semigroup on $X$ and let $B$ be a closed operator with $\mathcal{D}(B)\supseteq\mathcal{D}(A)$. Show that the operator defined by $\mathcal{A}(u,v)=(Bv,Av)$ generates an analytic semigroup on $X\times X$.

\item\label{ex:analytic-semigroup-strongly-damped-wave-equation} Discuss how analytic semigroups can be applied to the equation $u_{tt}=\Delta u_t+\Delta u$ with Dirichlet boundary conditions.

\end{enumerate}
